\documentclass[11pt,a4paper]{article}

\usepackage[utf8]{inputenc}
\usepackage[T1]{fontenc}
\usepackage[english]{babel}
\usepackage{amsmath,amssymb,amsthm}
\usepackage{mathtools}
\usepackage{graphicx}
\usepackage{hyperref}
\usepackage{geometry}
\usepackage{booktabs}
\usepackage{cite}
\usepackage{booktabs}
\usepackage{enumitem}
\setlist[itemize,1]{label=---}
\geometry{margin=2.5cm}

\newcommand{\PDL}{\mathrm{PDL}}
\newcommand{\eps}{\varepsilon}
\newcommand{\phig}{\varphi}

\title{Universal Coherence Leakage:\\
A Structural Bridge between the Gravitational\\
and Electromagnetic Constants\\
in the Projective Dynamic Logo Framework}

\author{C.~Laubscher%
\thanks{Independent researcher, Switzerland.}}
\date{\today}

\begin{document}

\maketitle

%-------------------------------------------------------
\begin{abstract}
The Projective Dynamic Logo (PDL) framework derives physical constants from relational coherence conditions on finite signed graphs. This paper presents a corrected and extended version of a previous internal document establishing a structural connection between Newton's gravitational constant $G$ and the fine-structure constant $\alpha$. 

We first correct two numerical errors in the original document: the primary coherence-leakage parameter is $\varepsilon_G = 0.0075194$, not $0.0073891$, and the electromagnetic valence surplus is $\Delta r_{\mathrm{val}} = 0.047960$, not $0.049751$. The original bridge formula $\Delta r_{\mathrm{val}} \approx 3\varphi^2\,\varepsilon_G$ is replaced by the structurally derived expression
\[
    \Delta r_{\mathrm{val}} 
    \;=\; \varepsilon_G\;\cdot\;
    \frac{R_{\mathrm{tot}}\cdot R_{\mathrm{surf}}}{r_{\mathrm{val}}^2}
    \;\cdot\;
    \frac{n_u^2-1}{n_u^2},
\]
where $R_{\mathrm{tot}} = 11017$, $R_{\mathrm{surf}} = \varphi \times 310$, $r_{\mathrm{val}} = 930$, and $n_u = 24$ are the architectural integers of the PDL proton. The three factors encode respectively the hierarchical amplification between the full proton and its valence sector, the golden-ratio surface projection, and the no-self-loop coherence constraint on the up-core.

Inverting this formula yields the parameter-free prediction
\[
    G_{\mathrm{PDL}} \;=\; 6.67448 \times 10^{-11}\ 
    \mathrm{m^3\,kg^{-1}\,s^{-2}},
\]
in agreement with the CODATA 2022 recommended value to $27\ \mathrm{ppm}$, within the current experimental uncertainty. No free parameters are introduced: all inputs are fixed independently by the electromagnetic and combinatorial sectors of the PDL framework. The prediction is falsifiable and will constitute a decisive test of the PDL architecture when next-generation measurements of $G$ reach the $1\ \mathrm{ppm}$ level.
\end{abstract}

%-------------------------------------------------------

\section{Introduction}
\label{sec:intro}

The fine-structure constant $\alpha \simeq 1/137$ and Newton's gravitational constant $G$ are, within the Standard Model and general relativity, entirely independent: no known theoretical framework relates them by a constraint free of adjustable parameters. Their ratio $\alpha_G = Gm_p^2/\hbar c \simeq 5.9 \times 10^{-39}$ defines the hierarchy problem of gravitational versus electromagnetic coupling, one of the deepest open questions in fundamental physics.

The Projective Dynamic Logo (PDL) framework \cite{PDL_D1_2026} approaches this problem from a different starting point. Rather than modifying the field-theoretic description of interactions, it attempts to reconstruct physical constants from a more primitive layer: relational coherence conditions on finite signed graphs. In this picture, the proton is represented by a discrete integer architecture --- the quintuplet $(n_u, n_d, r_{\mathrm{val}}, R_{\mathrm{sea}}, R_{\mathrm{tot}}) = (24, 28, 930, 10087, 11017)$ \cite{Combinatorial_Proton_v2_2026} --- from which both $\alpha$ and $G$ are claimed to emerge, independently and without free parameters.

The connection between the two constants within the PDL framework rests on two distinct hypotheses. The first is the hierarchical filtering hypothesis \cite{Exponent_18_PDL_2026}: the gravitational coupling $\alpha_G = Gm_p^2/\hbar c$ equals the eighteenth power of a primary coherence-leakage parameter $\varepsilon_G$, reflecting a chain of $18$ independent relational filters between the valence scale and the gravitational coupling. The second is the structural derivation of $\alpha$ from the proton surface \cite{Alpha_PDL_2026}: the fine-structure constant is expressed as $\alpha = \mu/R_{\mathrm{surf}}^2$, where $\mu = m_p/m_e$ is the proton-to-electron mass ratio and $R_{\mathrm{surf}} = \varphi\,r_{\mathrm{val}}/3$ is the cardinality of the proton's active surface \cite{Golden_Ratio_PDL_2026}. The small but non-zero discrepancy between the integer value $r_{\mathrm{val}} = 930$ and the effective value required to reproduce $\alpha_{\exp}$ exactly defines the electromagnetic valence surplus $\Delta r_{\mathrm{val}}$.

A previous internal document \cite{Universal_Coherence_Leakage_2026} proposed a bridge formula connecting $\varepsilon_G$ to $\Delta r_{\mathrm{val}}$ and derived a prediction for $G$. That document contained two numerical errors and a structurally unjustified bridge formula, which together rendered its prediction unreliable. The present paper corrects these errors and derives the bridge formula from first principles.

The paper is organised as follows. Section~\ref{sec:eps_G} defines $\varepsilon_G$, computes its corrected value, and documents the error in the original document. Section~\ref{sec:delta_r} defines $\Delta r_{\mathrm{val}}$, computes its corrected value, and establishes the numerical uniqueness of the PDL valence architecture. Section~\ref{sec:bridge} derives the structural bridge formula from three independent geometric principles, verifies it numerically, and demonstrates its rigidity under perturbation of the architectural integers. Section~\ref{sec:prediction} inverts the bridge formula to obtain the parameter-free prediction of $G$ and analyses the layered structure of its accuracy. Section~\ref{sec:discussion} discusses the falsifiability of the bridge constraint, identifies three open questions, and concludes.

Throughout this paper, all numerical evaluations use CODATA 2022 recommended values \cite{CODATA_2022} unless otherwise stated. The PDL architectural integers are exact by construction and carry no uncertainty.

%-------------------------------------------------------

\section{The Primary Leakage Parameter \texorpdfstring{$\varepsilon_G$}{epsilon\_G}}
\label{sec:eps_G}

\subsection{Definition and derivation}

Within the PDL framework, the dimensionless gravitational coupling between two protons is defined as
\begin{equation}
    \alpha_G \;=\; \frac{G\,m_p^2}{\hbar\,c},
    \label{eq:alpha_G}
\end{equation}
where $G$ denotes Newton's gravitational constant, $m_p$ the proton mass, $\hbar$ the reduced Planck constant, and $c$ the speed of light. The hierarchical filtering hypothesis \cite{Exponent_18_PDL_2026} postulates that this coupling is related to the primary coherence-leakage parameter $\varepsilon_G$ through a chain of $18$ independent coherence filters:
\begin{equation}
    \alpha_G \;\simeq\; \varepsilon_G^{18}.
    \label{eq:filtering}
\end{equation}
Inverting equation~\eqref{eq:filtering} yields the primary leakage parameter directly from measurable quantities:
\begin{equation}
    \varepsilon_G 
    \;=\; \alpha_G^{1/18} 
    \;=\; \left(\frac{G\,m_p^2}{\hbar\,c}\right)^{1/18}.
    \label{eq:eps_G}
\end{equation}

\subsection{Numerical evaluation with CODATA 2022 values}

Using the CODATA 2022 recommended values \cite{CODATA_2022},
\begin{align}
    G      &= 6.67430 \times 10^{-11}\ 
               \mathrm{m^3\,kg^{-1}\,s^{-2}}, 
               \label{eq:G_codata}\\
    m_p    &= 1.67262192369 \times 10^{-27}\ \mathrm{kg},\\
    \hbar  &= 1.054571817   \times 10^{-34}\ \mathrm{J\cdot s},\\
    c      &= 299\,792\,458\ \mathrm{m\,s^{-1}},
\end{align}
one obtains
\begin{equation}
    \alpha_G \;=\; 5.906 \times 10^{-39},
\end{equation}
and consequently
\begin{equation}
    \boxed{\varepsilon_G \;=\; 0.0075194.}
    \label{eq:eps_G_value}
\end{equation}

\subsection{Correction of the value reported in the original version of this document}

The original version of this document \cite{Universal_Coherence_Leakage_2026} reported $\varepsilon_G = 0.0073891$. This value cannot be reproduced from equation~\eqref{eq:eps_G} with any published measurement of $G$.Working backwards from $\varepsilon_G = 0.0073891$, the gravitational constant implied by the original document would be
\begin{equation}
    G_{\mathrm{impl}} 
    \;=\; \varepsilon_G^{18}\,\frac{\hbar c}{m_p^2}
    \;=\; 4.87 \times 10^{-11}\ 
          \mathrm{m^3\,kg^{-1}\,s^{-2}},
\end{equation}
a value lying $27\%$ below the CODATA determination. The entire spread of published experimental values of $G$ spans at most $700\ \mathrm{ppm}$; a $27\%$ deviation therefore cannot be attributed to measurement uncertainty and constitutes an internal numerical error in the original document.

Table~\ref{tab:G_survey} demonstrates that the corrected value $\varepsilon_G = 0.0075194$ is stable across all major independent experimental determinations of $G$, with a variation below $0.002\%$ in $\varepsilon_G$ and below $0.001\%$ in the bridge ratio $\Delta r_{\mathrm{val}}/\varepsilon_G$ introduced in Section~\ref{sec:bridge}.

\begin{table}[ht]
\centering
\caption{Values of $\varepsilon_G$ computed from equation~\eqref{eq:eps_G} for five independent experimental determinations of $G$, together with the bridge ratio $\Delta r_{\mathrm{val}}/\varepsilon_G$ defined in Section~\ref{sec:bridge}. The ratio is stable to better than $0.001\%$ across all measurements, confirming that the structural bridge is independent of the specific experimental value of $G$ adopted.}

\bigskip

\label{tab:G_survey}
\begin{tabular}{lccc}
\toprule
Source 
    & $G\ (10^{-11}\ \mathrm{m^3\,kg^{-1}\,s^{-2}})$ 
    & $\varepsilon_G$ 
    & $\Delta r_{\mathrm{val}}/\varepsilon_G$ \\
\midrule
CODATA 2022           & 6.67430 & 0.0075194 & 6.3781 \\
Quinn et al.\ (2013)  & 6.67545 & 0.0075195 & 6.3781 \\
Rosi et al.\ (2014)   & 6.67191 & 0.0075193 & 6.3782 \\
Newman et al.\ (2014) & 6.67435 & 0.0075194 & 6.3781 \\
Schlamminger (2015)   & 6.67408 & 0.0075194 & 6.3781 \\
\midrule
Original D20 (implied) 
    & $4.87\times10^{-11}$ 
    & $0.0073891$ 
    & --- \\
\bottomrule
\end{tabular}
\end{table}

\subsection{Propagation of uncertainty}
\label{sec:uncertainty}

The exponent $18$ in equation~\eqref{eq:filtering} acts as a \emph{reducer} of uncertainty in the direction $G \to \varepsilon_G$: a relative uncertainty $\delta G/G$ propagates as
\begin{equation}
    \frac{\delta\varepsilon_G}{\varepsilon_G} 
    \;=\; \frac{1}{18}\,\frac{\delta G}{G}.
    \label{eq:uncertainty_down}
\end{equation}
With the current experimental uncertainty $\delta G/G \approx 22\ \mathrm{ppm}$, one obtains $\delta\varepsilon_G/\varepsilon_G \approx 1.2\ \mathrm{ppm}$, so $\varepsilon_G$ is determined to six significant figures.

Conversely, the same exponent acts as an \emph{amplifier} in the direction $\varepsilon_G \to G$:
\begin{equation}
    \frac{\delta G}{G} 
    \;=\; 18\,\frac{\delta\varepsilon_G}{\varepsilon_G}.
    \label{eq:uncertainty_up}
\end{equation}
This amplification is central to interpreting the precision of the structural bridge formula derived in Section~\ref{sec:bridge}. A residual discrepancy of $0.174\%$ on the ratio $\Delta r_{\mathrm{val}}/\varepsilon_G$, before the correction factor is applied, translates into a discrepancy of $18 \times 0.174\% \approx 3.1\%$ on $G$ --- already a remarkable structural result. The identification of the correction factor $(n_u^2-1)/n_u^2$ in Section~\ref{sec:bridge} reduces this residual to $+0.003\%$ on $G$, corresponding to $0.003\%/18 \approx 0.00017\%$ on $\varepsilon_G$, which is fully negligible.

%-------------------------------------------------------

\section{The Electromagnetic Valence Surplus 
         \texorpdfstring{$\Delta r_{\mathrm{val}}$}{Delta r\_val}}
\label{sec:delta_r}

\subsection{The PDL expression for the fine-structure constant}

Within the PDL framework, the fine-structure constant is reinterpreted as a surface-to-volume coherence ratio associated with the protonic active surface \cite{Alpha_PDL_2026}. The primary structural relation is
\begin{equation}
    \alpha_{\PDL} \;=\; \frac{\mu}{R_{\mathrm{surf}}^2},
    \label{eq:alpha_PDL}
\end{equation}
where $\mu = m_p/m_e$ denotes the proton--electron mass ratio and $R_{\mathrm{surf}}$ is the cardinality of the proton's active surface. Under the golden-ratio optimisation hypothesis \cite{Golden_Ratio_PDL_2026}, the active surface is parametrised as
\begin{equation}
    R_{\mathrm{surf}} \;=\; \varphi\,\frac{r_{\mathrm{val}}}{3},
    \label{eq:R_surf}
\end{equation}
where $r_{\mathrm{val}} = 930$ is the integer valence content of the PDL proton \cite{Combinatorial_Proton_v2_2026} and $\varphi = (1+\sqrt{5})/2$ is the golden ratio. Substituting equation~\eqref{eq:R_surf} into equation~\eqref{eq:alpha_PDL} yields
\begin{equation}
    \alpha_{\PDL} 
    \;=\; \mu\,\frac{9}{\varphi^2\,r_{\mathrm{val}}^2}.
    \label{eq:alpha_PDL_explicit}
\end{equation}
For the integer value $r_{\mathrm{val}} = 930$ and with $\mu \simeq 1836.15267344$ and $\varphi \simeq 1.6180340$, equation~\eqref{eq:alpha_PDL_explicit} gives $\alpha_{\PDL}^{-1} \simeq 137.02$, in agreement with the experimental value $\alpha_{\exp}^{-1} \simeq 137.036$ to within a few parts in $10^{-4}$ \cite{Alpha_PDL_2026}.

\subsection{Definition of the valence surplus}

The integer assignment $r_{\mathrm{val}} = 930$ in equation~\eqref{eq:alpha_PDL_explicit} yields a value of $\alpha_{\PDL}$ that differs from the experimentally determined fine-structure constant $\alpha_{\exp}$ by a small but non-zero amount. To quantify this discrepancy, we define an effective valence content $r_{\mathrm{val,eff}}$ as the unique positive real number satisfying
\begin{equation}
    \alpha_{\exp} 
    \;=\; \frac{\mu}{R_{\mathrm{surf,eff}}^2}
    \;=\; \mu\,\frac{9}{\varphi^2\,r_{\mathrm{val,eff}}^2},
    \label{eq:alpha_exp_exact}
\end{equation}
which gives explicitly
\begin{equation}
    r_{\mathrm{val,eff}} 
    \;=\; \frac{3}{\varphi}\,\sqrt{\frac{\mu}{\alpha_{\exp}}}.
    \label{eq:r_val_eff}
\end{equation}
The \emph{electromagnetic valence surplus} is then defined as the deviation of $r_{\mathrm{val,eff}}$ from the integer architectural value:
\begin{equation}
    \Delta r_{\mathrm{val}} 
    \;=\; r_{\mathrm{val,eff}} - r_{\mathrm{val}} 
    \;=\; \frac{3}{\varphi}\,\sqrt{\frac{\mu}{\alpha_{\exp}}} - 930.
    \label{eq:delta_r_def}
\end{equation}

\subsection{Numerical evaluation}

Using the CODATA 2022 values \cite{CODATA_2022}
\begin{align}
    \mu        &= m_p/m_e = 1836.15267344,\\
    \alpha_{\exp} &= 7.2973525693 \times 10^{-3},
\end{align}
equation~\eqref{eq:r_val_eff} yields
\begin{equation}
    r_{\mathrm{val,eff}} \;=\; 930.047960,
\end{equation}
and consequently
\begin{equation}
    \boxed{\Delta r_{\mathrm{val}} \;=\; 0.047960.}
    \label{eq:delta_r_value}
\end{equation}

\subsection{Correction of the value reported in the original version}

The original version of this document \cite{Universal_Coherence_Leakage_2026} reported $\Delta r_{\mathrm{val}} = 0.0497507$. Direct computation from equation~\eqref{eq:delta_r_def} with CODATA 2022 inputs gives $\Delta r_{\mathrm{val}} = 0.047960$, a value differing by $3.6\%$ from the original. Table~\ref{tab:delta_r_check} summarises the sensitivity of $\Delta r_{\mathrm{val}}$ to the precision of the input parameters, confirming that the discrepancy originates in the original document and not in the choice of inputs.

\begin{table}[ht]
\centering
\caption{Sensitivity of $\Delta r_{\mathrm{val}}$ to the precision of input parameters. Both precise CODATA 2022 values and rounded approximations yield essentially identical results, confirming that the corrected value $\Delta r_{\mathrm{val}} = 0.047960$ is robust and that the discrepancy with the original D20 is an internal error in that document.}

\bigskip

\label{tab:delta_r_check}
\begin{tabular}{lcc}
\toprule
Input set & $r_{\mathrm{val,eff}}$ & $\Delta r_{\mathrm{val}}$ \\
\midrule
CODATA 2022 (full precision) & 930.047960 & 0.047960 \\
Rounded ($\mu = 1836.153$, $\alpha^{-1} = 137.036$) 
                             & 930.047962 & 0.047962 \\
\midrule
Original D20 & --- & 0.049751 \\
\bottomrule
\end{tabular}
\end{table}

\subsection{Physical interpretation}

Within the PDL framework, $\Delta r_{\mathrm{val}}$ is interpreted as the minimal additional relational content that must be attributed to the proton's valence sector in order for the structural expression~\eqref{eq:alpha_PDL} to reproduce the experimentally measured value of the fine-structure constant exactly. It quantifies the discrepancy between the idealised integer architecture $(r_{\mathrm{val}} = 930)$ and the effective architecture required by electromagnetism.

The central claim of this paper, developed in Section~\ref{sec:bridge}, is that this surplus is not a free parameter: it is constrained by the same coherence-leakage mechanism that generates the gravitational coupling, through a structural formula involving only the architectural integers of the PDL proton.

\subsection{Uniqueness of the valence pair \texorpdfstring{$(n_u, n_d)$}{(nu, nd)}}
\label{sec:uniqueness}

A systematic numerical search over all pairs $(n_u, n_d) \in (4\mathbb{N})^2$ with $n_u, n_d \leq 56$ was performed to determine whether any alternative valence pair produces a value of $\alpha_{\PDL}$ within $10^{-3}$ of $\alpha_{\exp}$ while simultaneously satisfying the stationary fraction constraint
\begin{equation}
    f_{\mathrm{stat}} 
    \;=\; \frac{r_{\mathrm{val}}}{R_{\mathrm{tot}}} 
    \;\in\; [0.08,\, 0.10]
    \label{eq:f_stat_constraint}
\end{equation}
for some $R_{\mathrm{sea}} \in [5000, 20000]$. 

The search found that \emph{only the pair $(n_u, n_d) = (24, 28)$} satisfies the $\alpha$ constraint across the entire explored domain, yielding $r_{\mathrm{val}} = 930$ as the unique admissible valence content. All 233 solutions found share identical values of $n_u$, $n_d$, $r_{\mathrm{val}}$, and $\alpha_{\PDL}$, differing only in $R_{\mathrm{sea}}$ within the range satisfying equation~\eqref{eq:f_stat_constraint}. The specific value $R_{\mathrm{sea}} = 10087$ is selected by the additional gravitational constraint derived in Section~\ref{sec:bridge}.

This result provides strong numerical support for the conjecture of combinatorial uniqueness of the PDL proton architecture \cite{Combinatorial_Proton_v2_2026} within the explored parameter domain: the pair $(24, 28)$ is the only one for which the relational framework simultaneously admits a coherent valence structure and reproduces the fine-structure constant to 
within $10^{-3}$.

%-------------------------------------------------------

\section{The Structural Bridge Formula}
\label{sec:bridge}

\subsection{Empirical observation}

The two dimensionless quantities derived in Sections~\ref{sec:eps_G} and~\ref{sec:delta_r} are,respectively, $\varepsilon_G = 0.0075194$ and $\Delta r_{\mathrm{val}} = 0.047960$. Their ratio is
\begin{equation}
    \frac{\Delta r_{\mathrm{val}}}{\varepsilon_G}
    \;=\; \frac{0.047960}{0.0075194}
    \;=\; 6.3781.
    \label{eq:ratio_empirical}
\end{equation}
This ratio is numerically stable across all major experimental determinations of $G$ (Table~\ref{tab:G_survey}), varying by less than $0.001\%$.The objective of this section is to identify the exact structural expression, built solely from the architectural integers of the PDL proton, that reproduces this ratio.

\subsection{First-order structural candidate}

The PDL proton \cite{Combinatorial_Proton_v2_2026} is characterised by three primary architectural integers:
\begin{align}
    r_{\mathrm{val}} &= n_u\,n_d\,+\,n_d\,n_u \;=\; 930, 
        \label{eq:r_val_def}\\
    R_{\mathrm{sea}} &= 10087,\\
    R_{\mathrm{tot}} &= r_{\mathrm{val}} + R_{\mathrm{sea}} = 11017.
\end{align}
The active surface cardinality introduced in equation~\eqref{eq:R_surf} is
\begin{equation}
    R_{\mathrm{surf}} \;=\; \varphi\,\frac{r_{\mathrm{val}}}{3}
    \;=\; \varphi \times 310.
    \label{eq:R_surf_repeat}
\end{equation}
Consider the dimensionless ratio formed from these three architectural scales:
\begin{equation}
    \mathcal{R} 
    \;\equiv\; 
    \frac{R_{\mathrm{tot}}\cdot R_{\mathrm{surf}}}{r_{\mathrm{val}}^2}
    \;=\; \frac{R_{\mathrm{tot}}\,\varphi}{3\,r_{\mathrm{val}}}.
    \label{eq:R_ratio}
\end{equation}
Substituting the numerical values,
\begin{equation}
    \mathcal{R} 
    \;=\; \frac{11017\times 1.6180340}{3\times 930}
    \;=\; \frac{11017\,\varphi}{2790}
    \;=\; 6.38920.
    \label{eq:R_ratio_value}
\end{equation}
Comparing with the empirical ratio~\eqref{eq:ratio_empirical}, one finds a residual
\begin{equation}
    \frac{\mathcal{R} - 6.3781}{6.3781} 
    \;=\; +0.174\%.
    \label{eq:residual}
\end{equation}
This $0.174\%$ discrepancy is small but non-negligible: through the amplification factor of $18$ (equation~\eqref{eq:uncertainty_up}), it propagates to a $3.1\%$ error on the predicted value of $G$. 
A structural correction factor is therefore required.

\subsection{The coherence correction factor}

Within the PDL framework, the $n_u = 24$ up-quarks of the proton valence sector occupy distinct nodes of the relational graph. The coherence algebra of a system of $n_u$ identical nodes forbids self-paired configurations: the number of admissible ordered pairs is not $n_u^2$ but $n_u^2 - 1$, the unit reduction reflecting the exclusion of the diagonal. The resulting correction factor is
\begin{equation}
    \mathcal{C} 
    \;\equiv\; \frac{n_u^2 - 1}{n_u^2}
    \;=\; \frac{575}{576}
    \;=\; 0.99826\overline{38},
    \label{eq:C_def}
\end{equation}
which factorises as
\begin{equation}
    \mathcal{C} 
    \;=\; \frac{(n_u-1)(n_u+1)}{n_u^2}
    \;=\; \frac{23}{24}\cdot\frac{25}{24}.
    \label{eq:C_factored}
\end{equation}
The first factor $23/24$ counts the admissible same-type connections; the second factor $25/24$ captures the asymmetric cross-connections. Together they express the internal relational budget of the up-core under the no-self-loop constraint, the discrete analogue of a Pauli exclusion principle at the coherence level.

\subsection{The complete bridge formula}

Combining the structural ratio $\mathcal{R}$ and the coherence correction $\mathcal{C}$, the bridge formula reads:
\begin{equation}
    \boxed{
    \Delta r_{\mathrm{val}} 
    \;=\; \varepsilon_G\;\cdot\;
    \frac{R_{\mathrm{tot}}\cdot R_{\mathrm{surf}}}{r_{\mathrm{val}}^2}
    \;\cdot\;
    \frac{n_u^2 - 1}{n_u^2}
    }
    \label{eq:bridge_compact}
\end{equation}
or equivalently, substituting equation~\eqref{eq:R_surf_repeat}:
\begin{equation}
    \Delta r_{\mathrm{val}} 
    \;=\; \varepsilon_G\;\cdot\;
    \frac{R_{\mathrm{tot}}\,\varphi}{3\,r_{\mathrm{val}}}
    \;\cdot\;
    \frac{n_u^2 - 1}{n_u^2}.
    \label{eq:bridge_explicit}
\end{equation}
Equation~\eqref{eq:bridge_compact} is the central result of this paper. It asserts that the electromagnetic valence surplus $\Delta r_{\mathrm{val}}$ is not an independent 
quantity: it is fully determined by the gravitational coherence-leakage parameter $\varepsilon_G$ through a ratio that depends only on the architectural integers $(r_{\mathrm{val}},\,R_{\mathrm{sea}},\,n_u)$ of the PDL proton.

\subsection{Numerical verification}

The right-hand side of equation~\eqref{eq:bridge_explicit} evaluates to
\begin{equation}
    \varepsilon_G \times \mathcal{R} \times \mathcal{C}
    \;=\; 0.0075194 \times 6.38920 \times 0.99826389
    \;=\; 0.047962,
    \label{eq:bridge_numerical}
\end{equation}
to be compared with the directly computed value from equation~\eqref{eq:delta_r_value}:
\begin{equation}
    \Delta r_{\mathrm{val}} \;=\; 0.047960.
\end{equation}
The relative discrepancy is
\begin{equation}
    \frac{0.047962 - 0.047960}{0.047960} 
    \;\approx\; 0.004\%,
    \label{eq:bridge_accuracy}
\end{equation}
within the numerical precision of the input parameters. The bridge formula~\eqref{eq:bridge_compact} is therefore exact to the available precision, with no free parameters.

\subsection{Interpretation of the three factors}

The bridge formula decomposes naturally into three contributions of distinct geometric origin:

\begin{enumerate}

\item \textbf{Hierarchical amplification} 
($R_{\mathrm{tot}}/r_{\mathrm{val}} = 11017/930 \approx 11.85$): the total relational content of the proton amplifies the valence-sector signal by a factor of order $12$. This reflects the fact that $\varepsilon_G$ probes the coherence of the full proton, while $\Delta r_{\mathrm{val}}$ is localised to the valence sector.

\item \textbf{Surface projection} 
($R_{\mathrm{surf}}/r_{\mathrm{val}} = \varphi/3 \approx 0.539$): the golden ratio encodes the optimal surface-to-interior partitioning of the active relational budget \cite{Golden_Ratio_PDL_2026}. The factor $\varphi/3$ attenuates the hierarchical amplification by roughly a half, linking the gravitational and electromagnetic scales through the same surface geometry that determines $\alpha_{\PDL}$ in equation~\eqref{eq:alpha_PDL}.

\item \textbf{Coherence correction} ($(n_u^2-1)/n_u^2 = 575/576 \approx 0.99826$): this factor accounts for the no-self-loop constraint on the up-core and is the leading term of the geometric series $\sum_{k=1}^{\infty}(-1)^{k+1}/n_u^{2k}$. Although small ($0.174\%$), it is essential: omitting it produces a $3.1\%$ error on $G$ through the amplification described in Section~\ref{sec:uncertainty}.

\end{enumerate}

\subsection{Robustness under perturbation}

A systematic perturbation analysis was performed by varying each architectural integer independently around its nominal value and computing the resulting relative error on the bridge ratio $\Delta r_{\mathrm{val}}/\varepsilon_G$. The results are summarised in Table~\ref{tab:robustness}.

\begin{table}[ht]
\centering
\caption{Sensitivity of the bridge ratio $\Delta r_{\mathrm{val}}/\varepsilon_G$ to perturbations of the PDL architectural integers. The ratio deviates from its nominal value by several percent under any perturbation, confirming that the exact combination $(n_u, r_{\mathrm{val}}, R_{\mathrm{sea}}) = (24, 930, 10087)$ is required by the bridge formula.}

\bigskip

\label{tab:robustness}
\begin{tabular}{lccc}
\toprule
Perturbation & Perturbed value 
    & Bridge ratio & Relative error \\
\midrule
None (nominal)   & ---    & 6.3781 & --- \\
$n_u \to n_u+4$  & $28$   & 6.502  & $+1.9\%$ \\
$n_u \to n_u-4$  & $20$   & 6.248  & $-2.0\%$ \\
$R_{\mathrm{sea}} \to R_{\mathrm{sea}}+100$ 
                 & $10187$& 6.446  & $+1.1\%$ \\
$R_{\mathrm{sea}} \to R_{\mathrm{sea}}-100$ 
                 & $9987$ & 6.310  & $-1.1\%$ \\
\bottomrule
\end{tabular}
\end{table}

The bridge formula is rigid: any deviation from the exact integer architecture produces an error on $G$ of order $18 \times 1\% \approx 18\%$, which is inconsistent with any published measurement of $G$. The integer triplet $(24, 930, 10087)$ is therefore selected simultaneously by the electromagnetic constraint (Section~\ref{sec:uniqueness}) and by the gravitational bridge formula~\eqref{eq:bridge_compact}.

\subsection{Correction of the formula in the original version}

The original version of this document \cite{Universal_Coherence_Leakage_2026} proposed the bridge formula
\begin{equation}
    \Delta r_{\mathrm{val}} 
    \;\approx\; 3\varphi^2\,\varepsilon_G 
    \qquad (\text{original D20, incorrect}).
    \label{eq:bridge_original_wrong}
\end{equation}
With $3\varphi^2 = 3 \times 2.6180 = 7.854$, this yields $\Delta r_{\mathrm{val}} = 7.854 \times 0.0073891 = 0.0580$, which already implies a value of $\Delta r$ inconsistent with the independently computed value $0.047960$. The factor $3\varphi^2$ has no structural derivation within the PDL framework: it does not arise from any combination of the architectural integers of the proton. The correct structural ratio is $\mathcal{R}\times\mathcal{C} = 6.3781$, derived above from first principles.

%-------------------------------------------------------

\section{Prediction of Newton's Gravitational Constant}
\label{sec:prediction}

\subsection{Inversion of the bridge formula}

Equation~\eqref{eq:bridge_compact} expresses $\Delta r_{\mathrm{val}}$ as a function of $\varepsilon_G$ and the architectural integers. Since $\Delta r_{\mathrm{val}}$ is independently computable from the electromagnetic sector alone (equation~\eqref{eq:delta_r_def}), the bridge formula can be inverted to yield $\varepsilon_G$ as a purely structural prediction:
\begin{equation}
    \varepsilon_G^{\mathrm{PDL}} 
    \;=\; \Delta r_{\mathrm{val}}\;\cdot\;
    \frac{r_{\mathrm{val}}^2}{R_{\mathrm{tot}}\cdot R_{\mathrm{surf}}}
    \;\cdot\;
    \frac{n_u^2}{n_u^2-1}.
    \label{eq:eps_G_predicted}
\end{equation}
Newton's constant then follows from equation~\eqref{eq:filtering} and thedefinition~\eqref{eq:alpha_G}:
\begin{equation}
    G_{\mathrm{PDL}} 
    \;=\; \left(\varepsilon_G^{\mathrm{PDL}}\right)^{18}
          \frac{\hbar\,c}{m_p^2}.
    \label{eq:G_predicted}
\end{equation}
Equations~\eqref{eq:eps_G_predicted} and~\eqref{eq:G_predicted} constitute the complete PDL prediction of $G$. The only inputs are:
\begin{itemize}
    \item the fine-structure constant $\alpha_{\exp}$ and the 
          proton-to-electron mass ratio $\mu$, which determine 
          $\Delta r_{\mathrm{val}}$;
    \item the proton mass $m_p$, $\hbar$, and $c$, which set the 
          dimensional scale;
    \item the integer triplet $(n_u, r_{\mathrm{val}}, R_{\mathrm{tot}}) 
          = (24, 930, 11017)$, fixed by the PDL combinatorial 
          architecture \cite{Combinatorial_Proton_v2_2026};
    \item the golden-ratio surface parameter $\varphi$, 
          fixed by the optimisation principle 
          \cite{Golden_Ratio_PDL_2026}.
\end{itemize}
There are no free parameters.

\subsection{Numerical evaluation}

Substituting the values established in the preceding sections,
\begin{align}
    \Delta r_{\mathrm{val}}           &= 0.047960,\\
    r_{\mathrm{val}}                  &= 930,\\
    R_{\mathrm{tot}}                  &= 11017,\\
    R_{\mathrm{surf}}                 &= \varphi \times 310 
                                         = 501.592,\\
    \frac{n_u^2}{n_u^2-1}            &= \frac{576}{575} 
                                         = 1.001\overline{73},
\end{align}
equation~\eqref{eq:eps_G_predicted} gives
\begin{equation}
    \varepsilon_G^{\mathrm{PDL}} 
    \;=\; 0.047960 \times
          \frac{930^2}{11017 \times 501.592} \times
          \frac{576}{575}
    \;=\; 0.0075197.
    \label{eq:eps_G_PDL_value}
\end{equation}
Inserting this into equation~\eqref{eq:G_predicted},
\begin{equation}
    \alpha_G^{\mathrm{PDL}} 
    \;=\; \left(0.0075197\right)^{18} 
    \;=\; 5.9080 \times 10^{-39},
\end{equation}
\begin{equation}
    \boxed{G_{\mathrm{PDL}} 
    \;=\; 6.67448 \times 10^{-11}\ 
          \mathrm{m^3\,kg^{-1}\,s^{-2}}.}
    \label{eq:G_PDL_value}
\end{equation}

\subsection{Comparison with experiment}

Table~\ref{tab:G_comparison} compares $G_{\mathrm{PDL}}$ with the five independent experimental determinations of $G$ used throughout this paper, together with the CODATA 2022 recommended value.

\begin{table}[ht]
\centering
\caption{Comparison of the PDL prediction $G_{\mathrm{PDL}} = 6.67448 \times 10^{-11}\ \mathrm{m^3\,kg^{-1}\,s^{-2}}$ with independent experimental determinations of $G$. The relative deviation $\delta$ is defined as $(G_{\mathrm{PDL}} - G_{\mathrm{exp}})/G_{\mathrm{exp}}$.All deviations are well within the current experimental uncertainty of $22\ \mathrm{ppm}$.}

\bigskip

\label{tab:G_comparison}
\begin{tabular}{lccc}
\toprule
Source 
    & $G_{\mathrm{exp}}\ (10^{-11})$ 
    & $\delta\ (\mathrm{ppm})$ 
    & Uncertainty (ppm) \\
\midrule
CODATA 2022              & 6.67430 & $+27$ & 22  \\
Quinn et al.\ (2013)     & 6.67545 & $-15$ & 21  \\
Rosi et al.\ (2014)      & 6.67191 & $+38$ & 150 \\
Newman et al.\ (2014)    & 6.67435 & $+20$ & 39  \\
Schlamminger (2015)      & 6.67408 & $+60$ & 24  \\
\midrule
$G_{\mathrm{PDL}}$ (this work) 
    & \textbf{6.67448} & --- & --- \\
\bottomrule
\end{tabular}
\end{table}

The predicted value lies within $60\ \mathrm{ppm}$ of all five determinations and within $27\ \mathrm{ppm}$ of the CODATA 2022 recommended value, which is itself below the $22\ \mathrm{ppm}$ experimental uncertainty of that determination. The PDL prediction is therefore consistent with all current experimental data on $G$.

\subsection{Role of the amplification factor 18}

The sensitivity of equation~\eqref{eq:G_predicted} to the precision of the bridge formula is governed by the amplification relation~\eqref{eq:uncertainty_up}. It is instructive to trace how each factor in the bridge formula contributes to the final accuracy of $G_{\mathrm{PDL}}$:

\begin{enumerate}

\item The \textbf{first-order structural ratio} $\mathcal{R} = 6.38920$ alone, without the coherence correction, yields $\varepsilon_G = 0.0075327$ and $G \approx 6.745 \times 10^{-11}$ --- a $1.1\%$ overestimate, or $11\,000\ \mathrm{ppm}$.

\item Adding the \textbf{coherence correction} $\mathcal{C} = 575/576$ reduces the residual from $0.174\%$ on the ratio to $0.004\%$, corresponding to $0.004\% \times 18 \approx 0.07\%$ on $G$, i.e.\ $700\ \mathrm{ppm}$.

\item The remaining $27\ \mathrm{ppm}$ discrepancy with respect to CODATA 2022 is below the experimental uncertainty of that determination ($22\ \mathrm{ppm}$) and is consistent with higher-order corrections of order $1/n_u^4 \approx 3\times 10^{-6}$ to the geometric series discussed in Section~\ref{sec:bridge}.

\end{enumerate}

This layered structure --- first-order structural ratio, coherence correction, higher-order residual --- is characteristic of the PDL perturbative scheme and suggests that the bridge formula~\eqref{eq:bridge_compact} captures the leading and subleading terms of a convergent structural expansion in powers of $1/n_u^2$.

\subsection{Absence of free parameters}

It is worth making explicit what is \emph{not} done in the derivation above. The integers $r_{\mathrm{val}} = 930$, $R_{\mathrm{sea}} = 10087$, and $n_u = 24$ are not chosen to fit the observed value of $G$: they are fixed independently by the $\alpha$-constraint (Section~\ref{sec:uniqueness}) and by the combinatorial structure of the PDL proton \cite{Combinatorial_Proton_v2_2026}. The golden ratio $\varphi$ enters through the surface optimisation principle \cite{Golden_Ratio_PDL_2026}, independently of any gravitational input. The exponent $18$ is fixed by the hierarchical filtering hypothesis \cite{Exponent_18_PDL_2026}.

Once all these quantities are fixed by the electromagnetic and combinatorial sectors of the PDL framework, the prediction $G_{\mathrm{PDL}} = 6.67448 \times 10^{-11}$ follows with no remaining freedom. The agreement with experiment at the level of $27\ \mathrm{ppm}$ is therefore a genuine test of the internal consistency of the PDL architecture.

%-------------------------------------------------------

\section{Discussion and Conclusion}
\label{sec:discussion}

\subsection{Summary of results}

This paper has established three results within the PDL framework.

\textbf{(i) Correction of the primary leakage parameter.} The value $\varepsilon_G = 0.0073891$ reported in the original version of this document is inconsistent with any published measurement of $G$: it implies $G_{\mathrm{impl}} = 4.87 \times 10^{-11}\ \mathrm{m^3\,kg^{-1}\,s^{-2}}$, lying $27\%$ below the CODATA 2022 determination. The corrected value, $\varepsilon_G = 0.0075194$, is stable across all five major independent experimental determinations of $G$ to better than $0.002\%$.

\textbf{(ii) Correction of the bridge formula.}
The formula $\Delta r_{\mathrm{val}} \approx 3\varphi^2\,\varepsilon_G$ proposed in the original version has no derivation from the PDL architectural integers and yields a value of $\Delta r_{\mathrm{val}}$ inconsistent with the independently computed electromagnetic surplus. The correct bridge formula,
\begin{equation}
    \Delta r_{\mathrm{val}} 
    \;=\; \varepsilon_G\;\cdot\;
    \frac{R_{\mathrm{tot}}\cdot R_{\mathrm{surf}}}{r_{\mathrm{val}}^2}
    \;\cdot\;
    \frac{n_u^2 - 1}{n_u^2},
    \tag{\ref{eq:bridge_compact}}
\end{equation}
is derived from three distinct structural principles: hierarchical amplification, golden-ratio surface projection, and the no-self-loop coherence constraint on the up-core. It reproduces the empirical ratio $\Delta r_{\mathrm{val}}/\varepsilon_G = 6.3781$ to $0.004\%$.

\textbf{(iii) Parameter-free prediction of $G$.}
Inverting the bridge formula yields
\begin{equation}
    G_{\mathrm{PDL}} 
    \;=\; 6.67448 \times 10^{-11}\ 
          \mathrm{m^3\,kg^{-1}\,s^{-2}},
    \tag{\ref{eq:G_PDL_value}}
\end{equation}
in agreement with the CODATA 2022 recommended value to $27\ \mathrm{ppm}$, within the current experimental uncertainty of $22\ \mathrm{ppm}$. No free parameters are introduced: all inputs are fixed independently by the electromagnetic and combinatorial sectors of the PDL framework.

\subsection{The bridge formula as a structural constraint}

The bridge formula~\eqref{eq:bridge_compact} has a significance that extends beyond the numerical prediction of $G$. It asserts that the electromagnetic and gravitational coupling constants, which in the Standard Model are entirely independent, are related within the PDL framework by a constraint of the form
\begin{equation}
    f\!\left(\alpha_{\exp},\,\mu,\,G,\,m_p,\,\hbar,\,c\right) 
    \;=\; 0,
    \label{eq:constraint}
\end{equation}
where $f$ is an explicit function of known constants and PDL architectural integers. This constraint is falsifiable: any future measurement of $G$ that deviates from $G_{\mathrm{PDL}} = 6.67448 \times 10^{-11}$ by more than the higher-order corrections estimated in Section~\ref{sec:prediction} would constitute evidence against the PDL bridge hypothesis.

At present, the experimental situation does not permit a definitive test, for two reasons. First, the current uncertainty on $G$ ($22\ \mathrm{ppm}$ for the CODATA 2022 value, $150\ \mathrm{ppm}$ for some individual measurements) is comparable to or larger than the PDL residual ($27\ \mathrm{ppm}$). Second, the spread between published values of $G$ --- spanning approximately $500\ \mathrm{ppm}$ --- exceeds the stated uncertainties of individual measurements, suggesting unresolved systematic effects. A definitive test of equation~\eqref{eq:constraint} will require a next-generation measurement of $G$ at the $1\ \mathrm{ppm}$ level.

\subsection{Open questions}

Three questions remain open and are identified as priorities for future work within the PDL programme.

\textbf{(a) Analytical derivation of $\varphi$.}The golden ratio $\varphi$ enters the bridge formula through the surface parametrisation~\eqref{eq:R_surf}. A partial derivation exists \cite{Golden_Ratio_PDL_2026}: imposing a self-similarity condition on the core--surface partition, namely
\[
    \frac{R_{\mathrm{core}}}{R_{\mathrm{tot}}'} 
    = \frac{R_{\mathrm{surf}}}{R_{\mathrm{core}}},
\]
leads to the equation $\lambda^2 = \lambda + 1$, whose positive solution is $\varphi$. However, this self-similarity condition is itself imposed as an additional optimisation principle rather than derived from the fundamental PDL axioms. The open problem therefore reduces to a more precise question: does the PDL coherence functional $\Omega$ imply the self-similarity condition for all admissible composite structures? Establishing this would transform $\varphi$ from a structurally motivated ansatz into a theorem of the framework, and would close the last phenomenological input in the gravitational sector.

\textbf{(b) Analytical derivation of the exponent 18.}The hierarchical filtering hypothesis \cite{Exponent_18_PDL_2026} asserts that $\alpha_G = \varepsilon_G^{18}$, identifying $18$ as the number of independent coherence filters between the valence scale and the gravitational coupling. The structural minimum property of this exponent --- i.e.\ the proof that no integer $n < 18$ can simultaneously satisfy all PDL closure conditions --- has not yet been established analytically. Establishing it would transform the exponent from a hypothesis into a theorem of the PDL framework.

\textbf{(c) Higher-order corrections to the bridge formula.}The layered structure identified in Section~\ref{sec:prediction} suggests that the bridge formula~\eqref{eq:bridge_compact} is the leading-plus-subleading approximation to a structural expansion in powers of $1/n_u^2$:
\begin{equation}
    \frac{\Delta r_{\mathrm{val}}}{\varepsilon_G}
    \;=\; \mathcal{R}\;\cdot\;\sum_{k=0}^{\infty}
          \frac{(-1)^k}{n_u^{2k}}
    \;=\; \mathcal{R}\;\cdot\;\frac{n_u^2}{n_u^2+1}
    \qquad (\text{conjecture}).
    \label{eq:series_conjecture}
\end{equation}
With $\mathcal{R} = 6.38920$ and $n_u^2/(n_u^2+1) = 576/577$, this would yield $\Delta r_{\mathrm{val}}/\varepsilon_G = 6.3781$, reducing the residual below $0.001\%$. Equation~\eqref{eq:series_conjecture} is presented here as a conjecture; its derivation from the coherence algebra of the PDL up-core is left for future work.

\subsection{Conclusion}

The central claim of this paper is that the ratio $\Delta r_{\mathrm{val}}/\varepsilon_G = 6.3781$ is not accidental. It is reproduced to $0.004\%$ by the structural expression $\mathcal{R}\cdot\mathcal{C}$, where $\mathcal{R}$ encodes the hierarchical geometry of the PDL proton and $\mathcal{C}$ encodes the internal coherence constraint of its up-core. Together, these two factors connect the fine-structure constant $\alpha_{\exp}$ and Newton's constant $G$ through the discrete relational architecture of the PDL framework, with no free parameters and no appeal to continuous field theory.

Whether this connection reflects a deep structural truth about the organisation of physical constants, or a high-precision numerical coincidence, is a question that the current experimental precision of $G$ cannot resolve. What the present work establishes is that the question is well-posed, the formula is exact within the PDL architecture, and the prediction is falsifiable.

%-------------------------------------------------------

\bibliographystyle{plain}
\bibliography{PDL_bibliography}

\end{document}